\section{1994년 제2판의 Bulloch와 Raz의 편집자 주 (Editors'
Note)}\label{uxb144-uxc81c2uxd310uxc758-bullochuxc640-razuxc758-uxd3b8uxc9d1uxc790-uxc8fc-editors-note}

\textbf{(제2판을 위해 작성됨)}

{\small

\emph{The Concept of Law}는 출간된 지 불과 몇 년 만에 영어권 세계와 그
너머에서 법리학(jurisprudence)이 이해되고 연구되는 방식을 변모시켰다. 이
책의 막대한 영향은 이 책과 그 교설들(doctrines)을 논의하는 수많은
출판물들로 이어졌으며, 그것은 법이론(legal theory)의 맥락에서뿐 아니라
정치철학(political philosophy)과 도덕철학(moral philosophy)에서도
그러했다.

하트(Hart)는 여러 해 동안 \emph{The Concept of Law}에 한 장(chapter)을
추가할 생각을 염두에 두고 있었다. 그는 그토록 큰 영향을 미친 본문을
자잘하게 손대기를 원하지 않았고, 그의 뜻에 따라 본문은 사소한
정정들(minor corrections)을 제외하고는 변경 없이 여기 출판된다. 그러나
그는 이 책에 관한 많은 논의들에 응답하기를 원했다. 그는 자신의 입장을
잘못 이해한 이들에 맞서 그 입장을 방어하고, 근거 없는 비판을 논박하며,
그가 보기에도 동등하게 중요하게는 정당한 비판의 힘(force)을 인정하고,
그러한 지점들에 대응하기 위해 책의 교설들(doctrines)을 조정하는 방식들을
제안하고자 했다. 새로운 장은 처음에는 머리말(preface)로 생각되었으나
최종적으로는 후기(postscript)로 생각되었고, 그가 사망할 당시
미완성이었다. 이는 부분적으로만 그의 꼼꼼한 완벽주의(perfectionism)
때문이었다. 그것은 또한 이 기획의 현명함에 관한 지속적인 의문들, 그리고
원래 구상된 책의 테제들(theses)이 가진 생동력과 통찰을 자신이 충분히
살려낼 수 있을지에 관한 끈질긴 불확실성 때문이기도 했다. 그럼에도 그는
많은 중단에도 불구하고 후기 작업을 계속했고, 사망 당시 의도된 두 절들
가운데 첫째 절은 거의 완성되어 있었다.

제니퍼 하트(Jennifer Hart)가 우리에게 초고들을 살펴보고 그 안에 출판
가능한 것이 있는지 결정해 달라고 요청했을 때, 우리의 가장 우선적인
생각은 하트가 만족하지 않았을 어떤 것도 출판되도록 해서는 안 된다는
것이었다. 그러므로 우리는 후기의 첫째 절 대부분이 그처럼 완성된 상태에
있음을 발견하고 기뻤다. 우리는 둘째 절을 위한 손글씨 주석들만
발견했는데, 그것들은 출판 가능하기에는 너무 단편적이고
미정형(inchoate)이었다. 대조적으로 첫째 절은 여러 판본들(versions)로
존재했고, 타자되고, 수정되고, 다시 타자되고, 다시 수정되어 있었다. 가장
최근의 판본조차도 그가 최종 상태로 생각한 것은 분명 아니었다. 연필과
바이로펜(Biro)으로 된 수많은 변경사항들(alterations)이 있다. 더욱이
하트는 더 이른 판본들을 버리지 않았고, 손 가까이에 있던 어느 판본에든
계속 작업한 것으로 보인다. 이것이 편집 작업을 더 어렵게 만들었지만, 지난
2년 동안 도입된 변경들은 대부분 문체적 뉘앙스의 변경들이었고, 그 사실
자체가 그가 본문을 있는 그대로 본질적으로 만족스러워했다는 점을
나타냈다.

우리의 과제는 대안적 판본들(alternative versions)을 비교하고, 그것들이
일치하지 않을 때에는 그중 하나에만 나타나는 본문 조각들(segments of
text)이 다른 판본들에서 빠져 있는 이유가 그가 그것들을 버렸기 때문인지,
아니면 모든 수정사항들(emendations)을 편입한 하나의 판본을 그가 한 번도
갖지 못했기 때문인지를 확정하는 것이었다. 출판된 텍스트는 하트가 버리지
않았고 그가 계속 수정한 텍스트 판본들에 나타나는 모든 수정사항들을
포함한다. 때때로 텍스트 자체는 비정합적이었다. 종종 이것은 타자 담당자가
원고를 잘못 읽은 결과였음이 틀림없고, 하트는 그 오류들을 언제나
알아차리지는 못했다. 다른 때에는 그것이 틀림없이, 그가 끝내 살아서 하지
못한 최종 원고 작성 단계에서 정리되었을 법한, 작성 과정에서 문장들이
자연스럽게 꼬이는 방식 때문이었다. 이런 경우 우리는 원래 텍스트를
복원하거나 최소한의 개입으로 하트의 생각(thought)을 되살리려 했다.
제6절(재량(discretion)에 관하여)은 하나의 특별한 문제를 제기했다. 우리는
그 첫 단락의 두 판본을 발견했는데, 하나는 그 지점에서 끝나는 사본 안에
있었고, 다른 하나는 그 절의 나머지를 포함하는 사본 안에 있었다. 잘린
판본(truncated version)은 그의 가장 최근 수정들 중 많은 것을 편입한 사본
안에 있었고, 그가 결코 버린 적이 없었으며, 후기의 일반적 논의와도
부합했으므로, 우리는 두 판본 모두 출판되도록 허용하기로 결정했다.
이어지지 않은 판본은 미주(endnote)에 나타난다.

하트는 주석들을, 대부분 참조들(references)인 그 주석들을, 타자하도록 한
적이 결코 없었다. 그는 손글씨 주석 판본을 가지고 있었고, 그
단서들(cues)은 본문 가장 이른 타자 사본에서 가장 쉽게 추적되었다. 나중에
그는 가끔 여백 논평에 참조들을 덧붙였지만, 대체로 이것들은 불완전했고,
때로는 그 참조를 추적할 필요만을 표시했다. 티머시 엔디컷(Timothy
Endicott)은 모든 참조들을 확인했고, 불완전한 것들을 모두 추적했으며,
하트가 원천(source)을 지목하지 않고 드워킨(Dworkin)을 인용하거나 그를
가깝게 풀어쓴 곳에는 참조들을 추가했다. 엔디컷은 또한 인용문들이
부정확했던 곳에서 본문을 바로잡았다. 폭넓은 조사와
수완(resourcefulness)을 수반한 이 작업 과정에서, 그는 위에서 제시된 편집
지침들에 부합하는 본문에 대한 몇 가지 정정사항들도 제안했고, 우리는 이를
감사히 편입하였다.

우리 마음속에는, 기회가 주어졌다면 하트가 출판 전에 텍스트를 더 다듬고
개선했을 것이라는 데 아무 의심도 없다. 그러나 우리는 출판된
후기(postscript)가 드워킨의 많은 논변들(arguments)에 대한 그의 숙고된
응답을 담고 있다고 믿는다.

}

\vspace{1em}

\begin{flushright}
\textbf{페넬로페 A. 불로흐 (Penelope A. Bulloch) 조셉 라즈 (Joseph Raz)}
\\ 1994
\end{flushright}
